% Source-derived chapter 05: Lefschetz classes
% Original source: standard-conjectures-abelian-fourfolds
\chapter{Lefschetz classes}
\label{standard-conjectures-abelian-fourfolds-chapter-05}
\source{standard-conjectures-abelian-fourfolds}

\begin{remark}[Source section]
\label{standard-conjectures-abelian-fourfolds-chapter-05-source}
\source{standard-conjectures-abelian-fourfolds}
\source{standard-conjectures-abelian-fourfolds}
This chapter reproduces the corresponding section of the retrieved
LaTeX source, including its definitions, statements, examples, and
proofs. It has not yet been translated into Lean.
\notready
\end{remark}

% The source section title was: Lefschetz classes
\label{sectionlefschetz}
\source{standard-conjectures-abelian-fourfolds}
The standard conjecture of Hodge type is known for divisors (Theorem \ref{segre}). In this section we explain how this implies the standard conjecture of Hodge type   for algebraic classes on abelian varieties that are linear combinations of intersections of divisors. This has been already pointed out by Milne \cite[Remark 3.7]{pola} using a different argument.

Throughout the section $A$ is a polarized abelian variety of dimension $g$. We will work with the motive $\mathfrak{h}^{n}(A)$ from Theorem \ref{classic} and the pairing  $\langle \cdot, \cdot \rangle_{n,\mot} $   from Definition \ref{defpairingmot}. 
%Recall also that the pairing $\langle \cdot, \cdot \rangle_n$ on $\mathcal{Z}^{n}_{\num}(A)_{\Q}$ from Definition  \ref{defpairing} is induced by $\langle \cdot, \cdot \rangle_{2n,\mot} $.
\begin{definition}
\source{standard-conjectures-abelian-fourfolds}
\notready\label{defexotic} Let  $\Q[\mathcal{Z}^{1}_{\num}(A)_{\Q}]$ be the subalgebra of $\mathcal{Z}^{*}_{\num}(A)_{\Q}$ generated by $\mathcal{Z}^{1}_{\num}(A)_{\Q}$.  An element of $\Q[\mathcal{Z}^{1}_{\num}(A)_{\Q}]$ is called a  Lefschetz class. The subspace  of Lefschetz classes in $\mathcal{Z}^{n}_{\num}(A)_{\Q}$ is denoted by  $\mathcal{L}^{n}(A)$.
\source{standard-conjectures-abelian-fourfolds}
%Define  $\mathcal{E}^{n}(A)$ as the subspace of  $\mathcal{Z}^{n}_{\num}(A)_{\Q}$ which is the orthogonal complement of $\mathcal{L}^{n}(A)$ with respect to the pairing $\langle \cdot, \cdot \rangle_n$. An element in $\mathcal{E}^{n}(A)$ is called an exotic class.
\end{definition}
\begin{remark}
\source{standard-conjectures-abelian-fourfolds}
\notready\label{nonsisa}
\source{standard-conjectures-abelian-fourfolds}
Given two positive integers $a,b$ such that $a\cdot b \leq g$ we have the canonical inclusions
\begin{equation} \label{eq:nonsisa}
\source{standard-conjectures-abelian-fourfolds}
\mathfrak{h}^{a\cdot b}(A)=\Sym^{a \cdot b} \mathfrak{h}^1(A) \subset (\Sym^{a} \mathfrak{h}^1(A))^{\otimes b}  = \mathfrak{h}^a(A)^{\otimes b}.
\end{equation}
  induced by Theorem \ref{classic}(\ref{kunn}). In particular any pairing on the right hand side induces a pairing on the left hand side. \end{remark}
\begin{lemma}
\source{standard-conjectures-abelian-fourfolds}
\notready\label{accoppiamenti1}
\source{standard-conjectures-abelian-fourfolds}
Fix an integer $n$ such that   $ 2 \leq 2n\leq g.$ Then the three motivic pairings $\langle \cdot, \cdot \rangle_{1,\mot}^{\otimes 2n}$,  $\langle \cdot, \cdot \rangle_{2,\mot}^{\otimes n}$ and $\langle \cdot, \cdot \rangle_{2n,\mot}$ on $\mathfrak{h}^{2n,\prim}(A)\in \NUM(k)_\Q$ coincide up to a positive rational scalar. Moreover, for each of these pairings, the Lefschetz decomposition $\mathfrak{h}^{2n}(A) = \mathfrak{h}^{2n,\prim}(A) \oplus \mathfrak{h}^{2n-2,\prim}(A) (-1)$ is orthogonal.\end{lemma}


\begin{proof}
 Our statement holds true even at homological level and not just numerically. For this, it is enough to check the statement after realization, for instance on the cohomology group $H_\ell^{2n,\prim}(A)$. 

To do this, take the moduli space of polarized abelian varieties (in mixed characteristic, with some fixed level structure). These pairings are defined on  the relative cohomology of the abelian scheme. Recall that the Zariski closure of the monodromy group associated to this relative cohomology is $\GSp_{2g}= \GSp(H_\ell^{1}(A))$ and that $H_\ell^{2n,\prim}(A)\subset  H_\ell^{2n }(A) = \Lambda^{2n}H_\ell^{1}(A)$  is a geometrically irreducible representation. This implies that these pairings coincide up to a scalar, by Schur's lemma. As these pairings are also defined on Betti cohomology this scalar must be rational. Moreover, they are polarizations of the underlying Hodge structure, so this scalar must be positive.

For the orthogonality part, the argument is the same. If, for a fixed pairing, the decomposition was not orthogonal, we would have a nonzero map between $H_\ell^{2n,\prim}(A)$ and $H_\ell^{2n-2}(A)(-1)^\vee$ which would be $\GSp_{2g}$-equivariant. This is impossible again by Schur's lemma.
 \end{proof}
\begin{proposition}
\source{standard-conjectures-abelian-fourfolds}
\notready\label{oklefschetz}
\source{standard-conjectures-abelian-fourfolds}
Fix an integer $n$ such that   $ 2 \leq 2n\leq g.$ The pairings $\langle \cdot, \cdot \rangle_{1,\mot }^{\otimes 2n}$,  $\langle \cdot, \cdot \rangle_{2,\mot}^{\otimes n}$ and $\langle \cdot, \cdot \rangle_{2n,\mot} $  on  $\mathcal{Z}^{n}_{\num}(A)_{\Q}$ are positive definite if and only if anyone of them is so. Moreover, they are positive definite on $\mathcal{L}^{n}(A)$.\end{proposition}
\begin{proof}
By Lemma \ref{accoppiamenti1} the Lefschetz decomposition \[\mathfrak{h}^{2n}(A) = \mathfrak{h}^{2n,\prim}(A) \oplus \mathfrak{h}^{2n-2,\prim}(A) (-1)\] is orthogonal with respect to any of these pairings, so, arguing by induction on $n$, it is enough to check positivity on algebraic classes of $ \mathfrak{h}^{2n,\prim}(A)$. Again by Lemma \ref{accoppiamenti1} the positivity on the primitive part does not depend on the pairing.



The argument just above works also for Lefschetz classes, indeed each component in the Lefschetz decomposition of a Lefschetz class is again a Lefschetz class  \cite[p. 640]{Milnelef}. Hence we can check positivity for one of the pairings, we will do it for $\langle \cdot, \cdot \rangle_{2,\mot}^{\otimes n}$. Notice that, by construction, the restriction of $\langle \cdot, \cdot \rangle_{2,\mot}^{\otimes n}$ to algebraic classes is $\langle \cdot, \cdot \rangle_{1}^{\otimes n}$, see Definition \ref{defpairing}.


Now, by (\ref{eq:nonsisa}), we have the inclusion
 \[\mathcal{Z}^n_{\num}(A)_{\Q}=\Hom_{\NUM(k)_\Q}(\one, \mathfrak{h}^{2n}(A)(n)) \subset  \Hom_{\NUM(k)_\Q}(\one,(\mathfrak{h}^2(A)(1))^{\otimes n}).\] 
 Moreover Theorem \ref{classic}(\ref{kunn})   implies the equality
 \[\mathcal{L}^{n}(A)_{\Q} = \mathcal{Z}^n_{\num}(A)_{\Q} \cap \mathcal{Z}^1_{\num}(A)_{\Q}^{\otimes n}\]
 where the intersection is taken inside  $\Hom_{\NUM(k)_\Q}(\one,(\mathfrak{h}^2(A)(1))^{\otimes n})$.

On the other hand, the pairing $\langle \cdot, \cdot \rangle_1$ on \[\mathcal{Z}^1_{\num}(A)_{\Q}= \Hom_{\NUM(k)_\Q}(\one, \mathfrak{h}^2(A)(1))\] 
is positive definite by Theorem \ref{segre}. Hence the pairing $\langle \cdot, \cdot \rangle_1^{\otimes n}$ on \[\mathcal{Z}^1_{\num}(A)^{\otimes n}_{\Q} = \Hom_{\NUM(k)_\Q}(\one, \mathfrak{h}^2(A)(1))^{\otimes n}\] is positive definite as well. In particular its restriction to $\mathcal{L}^{n}(A)_{\Q}$ will also be positive definite. \end{proof}
